\documentclass[12pt]{article}
\usepackage[letterpaper,margin=1in]{geometry}
\usepackage[T1]{fontenc}
\usepackage[utf8]{inputenc}
\usepackage{lmodern}
\usepackage{microtype}
\usepackage{setspace}
\usepackage{tabularx}
\usepackage{enumitem}
\usepackage[hidelinks,bookmarksopen=true,bookmarksnumbered=false]{hyperref}

\IfFileExists{filing-info.tex}{% Shared filing metadata for Civil/Filing documents.
%
% Keeps party names, contact information, and caption fragments here so updates
% flow through all filing materials.

\providecommand{\FilingCourtName}{UNITED STATES DISTRICT COURT}
\providecommand{\FilingDistrictName}{WESTERN DISTRICT OF TEXAS}
\providecommand{\FilingDivisionName}{AUSTIN DIVISION}
\providecommand{\FilingDivisionShort}{AUSTIN}
\providecommand{\FilingDistrictPartOne}{Western}
\providecommand{\FilingDistrictPartTwo}{Texas}
\providecommand{\FilingCivilActionNumber}{}
% Filing date shown on signature blocks and court forms.
% Defaults to the compilation date; replace with a fixed literal if needed.
\providecommand{\FilingDate}{June 12, 2026}

\providecommand{\FilingPlaintiffName}{Cory J. Bennett}
\providecommand{\FilingPlaintiffNameUpper}{CORY J. BENNETT}

\providecommand{\FilingPlaintiffHomeLineOne}{}      % Redacted for privacy; use if needed for court forms.
\providecommand{\FilingPlaintiffHomeCityStateZip}{} %
\providecommand{\FilingPlaintiffHomeAddress}{\FilingPlaintiffHomeLineOne, \FilingPlaintiffHomeCityStateZip}

% Use these for court contact/service addresses.
\providecommand{\FilingPlaintiffMailingLineOne}{6705 W Highway 290}
\providecommand{\FilingPlaintiffMailingLineTwo}{Ste 607 PMB \#268}
\providecommand{\FilingPlaintiffMailingCityStateZip}{Austin, TX 78735-8407}
\providecommand{\FilingPlaintiffMailingAddress}{\FilingPlaintiffMailingLineOne, \FilingPlaintiffMailingLineTwo, \FilingPlaintiffMailingCityStateZip}
\providecommand{\FilingPlaintiffTelephone}{(512) 630-5607}
\providecommand{\FilingPlaintiffFax}{None}
\providecommand{\FilingPlaintiffEmail}{csquaredbennett@gmail.com}

\providecommand{\FilingDefendantUniversity}{The University of Texas at Austin}
\providecommand{\FilingDefendantHsiao}{Emily Hsiao}
\providecommand{\FilingDefendantScott}{James G. Scott}
\providecommand{\FilingDefendantRagsdale}{Sally K. Ragsdale}
\providecommand{\FilingDefendantBartroff}{Jay Bartroff}
\providecommand{\FilingDefendantBradley}{Kelli Bradley}
\providecommand{\FilingDefendantGraves}{Jeff Graves}

\providecommand{\FilingDefendantsInline}{\FilingDefendantUniversity; \FilingDefendantHsiao; \FilingDefendantScott; \FilingDefendantRagsdale; \FilingDefendantBartroff; \FilingDefendantBradley; and \FilingDefendantGraves}
\providecommand{\FilingDefendantsInlineNoAnd}{\FilingDefendantUniversity; \FilingDefendantHsiao; \FilingDefendantScott; \FilingDefendantRagsdale; \FilingDefendantBartroff; \FilingDefendantBradley; \FilingDefendantGraves}
\providecommand{\FilingDefendantsInlineUpper}{THE UNIVERSITY OF TEXAS AT AUSTIN; EMILY HSIAO; JAMES G. SCOTT; SALLY K. RAGSDALE; JAY BARTROFF; KELLI BRADLEY; and JEFF GRAVES}

\providecommand{\FilingDefendantsCaptionLineOne}{\FilingDefendantUniversity;}
\providecommand{\FilingDefendantsCaptionLineTwo}{\FilingDefendantHsiao; \FilingDefendantScott; \FilingDefendantRagsdale;}
\providecommand{\FilingDefendantsCaptionLineThree}{\FilingDefendantBartroff; \FilingDefendantBradley; \FilingDefendantGraves}

\providecommand{\FilingDefendantsShortLineOne}{\FilingDefendantUniversity; \FilingDefendantHsiao; \FilingDefendantScott;}
\providecommand{\FilingDefendantsShortLineTwo}{\FilingDefendantRagsdale; \FilingDefendantBartroff; \FilingDefendantBradley; \FilingDefendantGraves}

\providecommand{\FilingDefendantsPermissionLineOne}{\FilingDefendantUniversity; \FilingDefendantHsiao;}
\providecommand{\FilingDefendantsPermissionLineTwo}{\FilingDefendantScott; \FilingDefendantRagsdale;}
\providecommand{\FilingDefendantsPermissionLineThree}{\FilingDefendantBartroff; \FilingDefendantBradley; \FilingDefendantGraves}
}{%
  \IfFileExists{../filing-info.tex}{% Shared filing metadata for Civil/Filing documents.
%
% Keeps party names, contact information, and caption fragments here so updates
% flow through all filing materials.

\providecommand{\FilingCourtName}{UNITED STATES DISTRICT COURT}
\providecommand{\FilingDistrictName}{WESTERN DISTRICT OF TEXAS}
\providecommand{\FilingDivisionName}{AUSTIN DIVISION}
\providecommand{\FilingDivisionShort}{AUSTIN}
\providecommand{\FilingDistrictPartOne}{Western}
\providecommand{\FilingDistrictPartTwo}{Texas}
\providecommand{\FilingCivilActionNumber}{}
% Filing date shown on signature blocks and court forms.
% Defaults to the compilation date; replace with a fixed literal if needed.
\providecommand{\FilingDate}{June 12, 2026}

\providecommand{\FilingPlaintiffName}{Cory J. Bennett}
\providecommand{\FilingPlaintiffNameUpper}{CORY J. BENNETT}

\providecommand{\FilingPlaintiffHomeLineOne}{}      % Redacted for privacy; use if needed for court forms.
\providecommand{\FilingPlaintiffHomeCityStateZip}{} %
\providecommand{\FilingPlaintiffHomeAddress}{\FilingPlaintiffHomeLineOne, \FilingPlaintiffHomeCityStateZip}

% Use these for court contact/service addresses.
\providecommand{\FilingPlaintiffMailingLineOne}{6705 W Highway 290}
\providecommand{\FilingPlaintiffMailingLineTwo}{Ste 607 PMB \#268}
\providecommand{\FilingPlaintiffMailingCityStateZip}{Austin, TX 78735-8407}
\providecommand{\FilingPlaintiffMailingAddress}{\FilingPlaintiffMailingLineOne, \FilingPlaintiffMailingLineTwo, \FilingPlaintiffMailingCityStateZip}
\providecommand{\FilingPlaintiffTelephone}{(512) 630-5607}
\providecommand{\FilingPlaintiffFax}{None}
\providecommand{\FilingPlaintiffEmail}{csquaredbennett@gmail.com}

\providecommand{\FilingDefendantUniversity}{The University of Texas at Austin}
\providecommand{\FilingDefendantHsiao}{Emily Hsiao}
\providecommand{\FilingDefendantScott}{James G. Scott}
\providecommand{\FilingDefendantRagsdale}{Sally K. Ragsdale}
\providecommand{\FilingDefendantBartroff}{Jay Bartroff}
\providecommand{\FilingDefendantBradley}{Kelli Bradley}
\providecommand{\FilingDefendantGraves}{Jeff Graves}

\providecommand{\FilingDefendantsInline}{\FilingDefendantUniversity; \FilingDefendantHsiao; \FilingDefendantScott; \FilingDefendantRagsdale; \FilingDefendantBartroff; \FilingDefendantBradley; and \FilingDefendantGraves}
\providecommand{\FilingDefendantsInlineNoAnd}{\FilingDefendantUniversity; \FilingDefendantHsiao; \FilingDefendantScott; \FilingDefendantRagsdale; \FilingDefendantBartroff; \FilingDefendantBradley; \FilingDefendantGraves}
\providecommand{\FilingDefendantsInlineUpper}{THE UNIVERSITY OF TEXAS AT AUSTIN; EMILY HSIAO; JAMES G. SCOTT; SALLY K. RAGSDALE; JAY BARTROFF; KELLI BRADLEY; and JEFF GRAVES}

\providecommand{\FilingDefendantsCaptionLineOne}{\FilingDefendantUniversity;}
\providecommand{\FilingDefendantsCaptionLineTwo}{\FilingDefendantHsiao; \FilingDefendantScott; \FilingDefendantRagsdale;}
\providecommand{\FilingDefendantsCaptionLineThree}{\FilingDefendantBartroff; \FilingDefendantBradley; \FilingDefendantGraves}

\providecommand{\FilingDefendantsShortLineOne}{\FilingDefendantUniversity; \FilingDefendantHsiao; \FilingDefendantScott;}
\providecommand{\FilingDefendantsShortLineTwo}{\FilingDefendantRagsdale; \FilingDefendantBartroff; \FilingDefendantBradley; \FilingDefendantGraves}

\providecommand{\FilingDefendantsPermissionLineOne}{\FilingDefendantUniversity; \FilingDefendantHsiao;}
\providecommand{\FilingDefendantsPermissionLineTwo}{\FilingDefendantScott; \FilingDefendantRagsdale;}
\providecommand{\FilingDefendantsPermissionLineThree}{\FilingDefendantBartroff; \FilingDefendantBradley; \FilingDefendantGraves}
}{% Shared filing metadata for Civil/Filing documents.
%
% Keeps party names, contact information, and caption fragments here so updates
% flow through all filing materials.

\providecommand{\FilingCourtName}{UNITED STATES DISTRICT COURT}
\providecommand{\FilingDistrictName}{WESTERN DISTRICT OF TEXAS}
\providecommand{\FilingDivisionName}{AUSTIN DIVISION}
\providecommand{\FilingDivisionShort}{AUSTIN}
\providecommand{\FilingDistrictPartOne}{Western}
\providecommand{\FilingDistrictPartTwo}{Texas}
\providecommand{\FilingCivilActionNumber}{}
% Filing date shown on signature blocks and court forms.
% Defaults to the compilation date; replace with a fixed literal if needed.
\providecommand{\FilingDate}{June 12, 2026}

\providecommand{\FilingPlaintiffName}{Cory J. Bennett}
\providecommand{\FilingPlaintiffNameUpper}{CORY J. BENNETT}

\providecommand{\FilingPlaintiffHomeLineOne}{}      % Redacted for privacy; use if needed for court forms.
\providecommand{\FilingPlaintiffHomeCityStateZip}{} %
\providecommand{\FilingPlaintiffHomeAddress}{\FilingPlaintiffHomeLineOne, \FilingPlaintiffHomeCityStateZip}

% Use these for court contact/service addresses.
\providecommand{\FilingPlaintiffMailingLineOne}{6705 W Highway 290}
\providecommand{\FilingPlaintiffMailingLineTwo}{Ste 607 PMB \#268}
\providecommand{\FilingPlaintiffMailingCityStateZip}{Austin, TX 78735-8407}
\providecommand{\FilingPlaintiffMailingAddress}{\FilingPlaintiffMailingLineOne, \FilingPlaintiffMailingLineTwo, \FilingPlaintiffMailingCityStateZip}
\providecommand{\FilingPlaintiffTelephone}{(512) 630-5607}
\providecommand{\FilingPlaintiffFax}{None}
\providecommand{\FilingPlaintiffEmail}{csquaredbennett@gmail.com}

\providecommand{\FilingDefendantUniversity}{The University of Texas at Austin}
\providecommand{\FilingDefendantHsiao}{Emily Hsiao}
\providecommand{\FilingDefendantScott}{James G. Scott}
\providecommand{\FilingDefendantRagsdale}{Sally K. Ragsdale}
\providecommand{\FilingDefendantBartroff}{Jay Bartroff}
\providecommand{\FilingDefendantBradley}{Kelli Bradley}
\providecommand{\FilingDefendantGraves}{Jeff Graves}

\providecommand{\FilingDefendantsInline}{\FilingDefendantUniversity; \FilingDefendantHsiao; \FilingDefendantScott; \FilingDefendantRagsdale; \FilingDefendantBartroff; \FilingDefendantBradley; and \FilingDefendantGraves}
\providecommand{\FilingDefendantsInlineNoAnd}{\FilingDefendantUniversity; \FilingDefendantHsiao; \FilingDefendantScott; \FilingDefendantRagsdale; \FilingDefendantBartroff; \FilingDefendantBradley; \FilingDefendantGraves}
\providecommand{\FilingDefendantsInlineUpper}{THE UNIVERSITY OF TEXAS AT AUSTIN; EMILY HSIAO; JAMES G. SCOTT; SALLY K. RAGSDALE; JAY BARTROFF; KELLI BRADLEY; and JEFF GRAVES}

\providecommand{\FilingDefendantsCaptionLineOne}{\FilingDefendantUniversity;}
\providecommand{\FilingDefendantsCaptionLineTwo}{\FilingDefendantHsiao; \FilingDefendantScott; \FilingDefendantRagsdale;}
\providecommand{\FilingDefendantsCaptionLineThree}{\FilingDefendantBartroff; \FilingDefendantBradley; \FilingDefendantGraves}

\providecommand{\FilingDefendantsShortLineOne}{\FilingDefendantUniversity; \FilingDefendantHsiao; \FilingDefendantScott;}
\providecommand{\FilingDefendantsShortLineTwo}{\FilingDefendantRagsdale; \FilingDefendantBartroff; \FilingDefendantBradley; \FilingDefendantGraves}

\providecommand{\FilingDefendantsPermissionLineOne}{\FilingDefendantUniversity; \FilingDefendantHsiao;}
\providecommand{\FilingDefendantsPermissionLineTwo}{\FilingDefendantScott; \FilingDefendantRagsdale;}
\providecommand{\FilingDefendantsPermissionLineThree}{\FilingDefendantBartroff; \FilingDefendantBradley; \FilingDefendantGraves}
}%
}

\renewcommand{\FilingCivilActionNumber}{1:26-cv-01601-RP-DH}

\setlength{\parindent}{0.5in}
\setlength{\parskip}{0pt}
\setlist[enumerate]{leftmargin=0.5in,itemsep=0.1\baselineskip,topsep=0.15\baselineskip,parsep=0pt,partopsep=0pt}
\doublespacing
\emergencystretch=3em
\pagestyle{plain}

\newcommand{\Caption}{%
  \begin{singlespace}
  \begin{center}
  \textbf{\FilingCourtName}\\
  \textbf{\FilingDistrictName}\\
  \textbf{\FilingDivisionName}
  \end{center}

  \vspace{0.45em}

  \noindent
  \begin{tabularx}{\textwidth}{@{}>{\raggedright\arraybackslash}X>{\raggedright\arraybackslash}p{2.95in}@{}}
  \textbf{\FilingPlaintiffNameUpper,} Plaintiff, & \textbf{Civil Action No.} \FilingCivilActionNumber \\
  \textbf{v.} & \\
  \textbf{\FilingDefendantsInlineUpper,} Defendants. & \\
  \end{tabularx}

  \vspace{1.0em}
  \end{singlespace}%
}

\begin{document}
\Caption

\begin{singlespace}
\begin{center}
\textbf{[PROPOSED] ORDER GRANTING PLAINTIFF'S}\\
\textbf{MOTION FOR PRELIMINARY INJUNCTION}
\end{center}
\end{singlespace}

Before the Court is Plaintiff Cory J. Bennett's Motion for Preliminary Injunction. The Court has considered the motion, declaration, appendix, response, reply, evidence presented at any hearing, and applicable law.

The Court finds that Plaintiff has shown a substantial likelihood of success on at least one claim under Title II of the Americans with Disabilities Act or Section 504 of the Rehabilitation Act; a substantial threat of losing Fall 2026 enrollment, course participation, and timely SDS coursework before final judgment; a balance of hardships favoring temporary relief; and a public interest in disability accommodations and enrollment decisions made before the Fall 2026 deadlines by officials who did not impose the October 2 restrictions, review Hsiao's revised SDS 320E lab instructions, dismiss Plaintiff's HOP complaint, decide Plaintiff's DIA appeal, or control the January D\&A coordinator assignment. If this Order issues after an ordinary registration, payment, or add deadline, UT can preserve Fall participation through an administrative add, a seat reservation when practicable, or a deadline extension.

Plaintiff intends to enroll in Fall 2026. Fall enrollment requires the University of Texas at Austin's Disability and Access office (``D\&A'') to issue and administer accommodations; the College of Natural Sciences (``CNS'') to prepare a current academic plan; and CNS or the Department of Statistics and Data Sciences (``SDS'') to decide the status of SDS 320E and the SDS 324E prerequisite. Bradley is the last Access Coordinator UT identified to Plaintiff. The October 2 SDS restrictions and related Student Conduct record remain unresolved. Page 1 of Exhibit 2 shows the Summer 2026 graduation timeline displaced by Spring non-enrollment, subject to Plaintiff's residence-hour petition. Page 2 shows an alternative sequence through Spring 2027 if SDS 320E must be repeated before SDS 324E.

IT IS ORDERED that Plaintiff's motion is GRANTED as follows:

\begin{enumerate}
\item \textbf{Excluded personnel.} Emily Hsiao, James G. Scott, Sally K. Ragsdale, Jay Bartroff, Michael Drew, Kelli Bradley, Emily Harrington, Jody Hughes, Patrick Joseph, Esteban Soto, Dawn Spinozza, and Jeff Graves shall not select, supervise, advise, draft, review, approve, veto, or decide Plaintiff's Fall 2026 accommodations, registration, course placement, SDS 320E status and SDS 324E prerequisite options, course participation, or review of decisions in those categories.

Within five business days, after reviewing records reasonably available to it, UT shall identify any additional employee whose records show that the employee did one or more of the following:

  \begin{enumerate}[label=(\alph*),leftmargin=0.45in]
  \item imposed, recommended, drafted, approved, or transmitted the September 30 or October 2, 2025 SDS and CNS restrictions, special communication requirements, instruction to contact the University of Texas Police Department, or recurring-behavior narrative;

  \item drafted, reviewed, approved, or confirmed Hsiao's revised SDS 320E lab instructions; consulted on those instructions; placed Plaintiff on Bradley's caseload because of that review; or controlled the January 2026 D\&A coordinator assignment or reassignment request;

  \item created, maintained, transmitted, or used the Student Conduct or Student Outreach and Support (``SOS'') record created from the SDS reports or later accommodation correspondence; or

  \item distributed Plaintiff's Department of Investigation and Adjudication (``DIA'') intake notes, deferred to D\&A, dismissed the HOP complaint, decided the DIA appeal, or responded to Plaintiff's request for review outside D\&A, DIA, Legal Affairs, the ADA Coordinator, SDS, or CNS.
  \end{enumerate}

Each additional employee identified under this paragraph shall also be excluded from the Fall 2026 decisions listed in this paragraph. Clerical handling, receipt of a copied communication, evidence preservation, or testimony alone does not require exclusion.

\item \textbf{Fall 2026 decisionmakers.} Within three business days, UT shall designate:

  \begin{enumerate}[label=(\alph*),leftmargin=0.45in]
  \item an Access Coordinator and supervising D\&A official;

  \item a CNS academic coordinator; and

  \item an academic official authorized to decide SDS prerequisites, substitutions, and equivalent coursework.
  \end{enumerate}

Each designee must be outside the personnel excluded under paragraph 1. Excluded personnel shall not select or supervise the designees, give them advice about Plaintiff's Fall 2026 decisions, review drafts of their decisions, approve those decisions, or veto them. UT Legal Affairs may represent UT in this litigation and preserve evidence. A Legal Affairs attorney excluded under paragraph 1 shall not serve as an administrative decisionmaker or control the designees' decisions about Plaintiff.

\item \textbf{Identification and certification.} Within five business days, UT shall file a notice that:

  \begin{enumerate}[label=(\alph*),leftmargin=0.45in]
  \item identifies the designees required by paragraph 2;

  \item identifies each additional employee excluded under paragraph 1; and

  \item certifies that each designee performed none of the acts listed in paragraph 1.
  \end{enumerate}

UT may file employee-role information under restricted access if necessary to protect confidential personnel information.

\item \textbf{Degree audit and academic plan.} Within five business days, the CNS academic coordinator shall provide Plaintiff with a current degree audit and written academic plan that separately identifies:

  \begin{enumerate}[label=(\alph*),leftmargin=0.45in]
  \item the remaining bachelor's degree requirements;

  \item the remaining Statistics and Data Science minor requirements;

  \item the certificate coursework Plaintiff has included in the Fall 2026 academic plan;

  \item the remaining residence-hour requirement and the official assigned to decide any pending or renewed residence-hour petition;

  \item elective credit that may be completed in Fall 2026 or a summer term;

  \item Fall 2026 courses with available seats or approved equivalents; and

  \item the earliest completion date under each available SDS 320E status and SDS 324E prerequisite option.
  \end{enumerate}

\item \textbf{Fall enrollment.} Within five business days, UT shall complete the administrative steps within its control needed for Plaintiff to register and enroll for Fall 2026. Plaintiff remains responsible for the academic prerequisites and financial obligations stated in the written academic plan.

If compliance occurs after a registration, payment, or add deadline, UT shall administratively add Plaintiff to an approved Fall 2026 course, reserve a seat when practicable, or extend the applicable deadline long enough for Plaintiff to complete registration.

\item \textbf{SDS 320E status and SDS 324E prerequisite determination.} Within seven business days, the academic official designated under paragraph 2 shall issue a written determination addressing:

  \begin{enumerate}[label=(\alph*),leftmargin=0.45in]
  \item whether SDS 320E will be repeated, substituted, restored as credit, or otherwise satisfied;

  \item whether Plaintiff may satisfy the SDS 324E prerequisite through an override, concurrent enrollment, equivalent course, independent study, or another academically equivalent arrangement; and

  \item the resulting course sequence and completion date under each available option.
  \end{enumerate}

An ordinary registration or add deadline that expires while the determination is pending shall not bar Plaintiff from enrolling in a course the official approves. UT shall reserve a seat when practicable, use an administrative add, or extend the applicable deadline. If the official declines every SDS 320E or SDS 324E option that would avoid an additional semester, the written determination shall identify the academic requirement affected by each option and explain why no equivalent arrangement is available.

\item \textbf{SDS instructional assignments.} UT shall place Plaintiff in available sections of SDS courses included in the written academic plan and taught, supervised, and graded by personnel not excluded under paragraph 1. When no such section is available, the CNS academic coordinator and academic official designated under paragraph 2 shall provide an academically equivalent instructional or grading arrangement if one is available under published academic standards. If no equivalent arrangement is available, they shall identify the affected requirement, the unavailable course or section, and the earliest available course or section that satisfies the same requirement.

\item \textbf{Fall accommodation plan.} Within seven business days, or earlier if needed to meet an applicable registration or add deadline, the D\&A designees shall confer with Plaintiff and issue a written Fall 2026 accommodation plan. The plan shall:

  \begin{enumerate}[label=(\alph*),leftmargin=0.45in]
  \item state the accommodations for sudden disability-related absence or incapacitation;

  \item permit Plaintiff to notify the designated contact after an absence when sudden disability-related incapacitation prevented earlier communication, and state how and when the missed work may be completed;

  \item identify the person responsible for each accommodation; and

  \item identify the official and deadline for deciding a disagreement over whether an approved accommodation was implemented.
  \end{enumerate}

\item \textbf{October 2 restrictions and Student Conduct record.} UT shall suspend the October 2, 2025 office-hours restriction, special communication requirements, and any instruction based on those restrictions that remains active. Plaintiff shall receive the instructor, teaching-assistant, office-hour, course-communication, make-up, and academic-support access available in his assigned courses, together with the accommodations stated in the Fall plan.

UT shall preserve every SDS, CNS, D\&A, Legal Affairs, ADA Coordinator, Student Conduct, SOS, DIA, URCS, and compliance record concerning the September 30 interaction, the October 2 restrictions, Hsiao's revised lab instructions, the HOP complaint and DIA appeal, the January D\&A coordinator assignment, and Fall 2026 implementation. Until further order, UT shall not use the disputed description of the September 30 interaction, the October 2 restrictions, or the unresolved Student Conduct record to deny or condition Plaintiff's enrollment, accommodations, course placement, prerequisite determination, instructor access, grading administration, student affiliation, or review rights.

\item \textbf{New restrictions and review.} If UT contends that a later interaction or other event supports a new restriction affecting Plaintiff, an official not excluded under paragraph 1 must decide whether to impose it. UT shall provide written notice of the factual basis and a prompt opportunity to respond. An immediate temporary measure shall receive written review by the D\&A and CNS designees within two business days.

Any review or appeal concerning Fall 2026 enrollment, accommodations, course placement, prerequisite determinations, or a disagreement over whether an accommodation was implemented shall be assigned to an official not excluded under paragraph 1. The reviewer shall receive the relevant record, permit Plaintiff to identify documents and witnesses, and issue a written decision with reasons.

\item \textbf{Administrative communications.} UT shall designate one coordinator for communications about Fall 2026 accommodations, advising, enrollment, and course planning. Plaintiff may bring a support person to a meeting about those topics. Written decisions shall identify the responsible official, the facts considered, the decision, and the review route.

Personnel excluded under paragraph 1 shall communicate with Plaintiff about his accommodations, enrollment, course assignments, prerequisite determination, or course participation only through litigation counsel, formal discovery, testimony, evidence-preservation requests, or the designated coordinator.

\item \textbf{Protection against retaliation.} UT and persons bound by Federal Rule of Civil Procedure 65(d)(2) shall not deny or alter Plaintiff's enrollment, accommodations, course assignment, prerequisite determination, instructor access, or implementation review because Plaintiff requested accommodations, objected to disability discrimination, requested records or witnesses, sought reassignment or recusal, filed agency complaints, or prosecuted this action.

\item \textbf{Compliance report.} Within ten business days, or earlier if needed to implement this Order before an applicable Fall 2026 registration or add deadline, UT shall file a report that identifies:

  \begin{enumerate}[label=(\alph*),leftmargin=0.45in]
  \item the Access Coordinator, supervising D\&A official, CNS academic coordinator, and academic official designated under paragraph 2;

  \item the certification required by paragraph 3;

  \item the registration steps completed;

  \item the Fall courses, seats, administrative adds, and equivalent arrangements offered;

  \item the SDS 320E status and SDS 324E prerequisite determination;

  \item the status of the Fall accommodation plan;

  \item the official assigned to any pending or renewed residence-hour petition; and

  \item each remaining academic prerequisite, financial obligation, or residence-hour requirement Plaintiff must satisfy.
  \end{enumerate}

\item \textbf{Security and duration.} Security under Rule 65(c) is waived. If the Court requires security, Plaintiff shall post security in the amount the Court sets. Plaintiff proposes \$100 as a nominal amount. This Order takes effect upon entry and remains in force until final judgment or further order. It binds UT Austin and the persons described in Federal Rule of Civil Procedure 65(d)(2) who receive actual notice of this Order.

\item \textbf{Issues reserved.} This Order does not decide damages, permanent record correction, permanent accommodation terms, course grades, degree conferral, or the ultimate merits of the claims and defenses.
\end{enumerate}

SIGNED on \underline{\hspace{1.4in}}, 2026.

\vspace{2.2em}
\noindent\rule{3.5in}{0.4pt}\\
UNITED STATES DISTRICT JUDGE

\end{document}
